% TOP_PROBLEM: {"id":30007713,"problem_number":"LOCAL-30007713","title":"Economic consequences of climate tipping points","category_id":21,"category":{"id":21,"name":"economics","display_name":"Economics","description":"Open economic theory statements, published research agendas, and proposed scoped specifications; question provenance and historical priority are qualified per record.","slug":"economics","order_index":21,"created_at":"2026-10-08T00:00:00Z"},"status":"provenance_pending","external_url":null,"legacy_ids":[],"published":false,"statement":"Bound the century-long welfare cost of one specified climate tipping event when migration, investment, and adaptation respond, separating physical uncertainty from economic uncertainty.","economics":{"original_id":202,"field":"Climate, environment and energy","kind":"empirical","formalization_basis":"proposed_specification","source_status":"not_verified","priority_status":"not_established","priority_note":"No readable primary passage explicitly stating this precise open question was verified. A supporting paper, abstract, or topic citation is insufficient to establish origin or unresolved status.","earliest_verified_reference":null,"current_reference":null,"formalization_note":"Proposed scoped research specification. The original catalogue prompt was a synthesis, and no canonical conjecture or verified original formulation is claimed. The supporting source, when retained, establishes context or existing results rather than this exact open question."},"statusQualification":"Provenance pending: an explicit published statement of this exact open question was not verified. This is a catalogue-authored candidate specification, not a certified open conjecture. Proposed scoped research specification. The original catalogue prompt was a synthesis, and no canonical conjecture or verified original formulation is claimed. The supporting source, when retained, establishes context or existing results rather than this exact open question.","statusReviewedAt":"2026-10-08"}
% FIRST_POSTED: 2026-10-08
% ENTRY_KIND: statement_only
% !TeX program = lualatex
\documentclass[11pt]{article}
\usepackage{fontspec}
\usepackage[a4paper,margin=25mm]{geometry}
\usepackage{amsmath,amssymb,mathtools}
\usepackage{xurl}
\usepackage[hidelinks]{hyperref}
\newcommand{\sref}[2]{\hyperlink{source:#1:#2}{[S#2]}}
\setlength{\emergencystretch}{3em}
\begin{document}
% problemId:30007713
\section*{Economic consequences of climate tipping points}\label{30007713}
\subsection{Definitions and mathematical statement}
Consider a calibrated world economy with a finite set of regions \(r\), annual periods \(t=0,\ldots,100\), and one specified physical tipping event, such as a major ice-sheet transition. Let \(T_t\) be global mean temperature, \(S_t\in\{0,1\}\) the irreversible event indicator, and \(q(T_t)\) its conditional transition probability when \(S_t=0\). Consumption \(C_{rt}>0\), adaptation investment \(A_{rt}\geq0\), capital, and migration are equilibrium outcomes. A stated damage map links the event to sea-level rise and regional production. For discount factor \(\beta\in(0,1)\), positive regional welfare weights \(w_r\), and increasing concave utility \(u_r\), define the incremental expected welfare loss
\[L=E\!\left[\sum_{t=0}^{100}\beta^t\sum_r w_r\{u_r(C_{rt}^{0})-u_r(C_{rt}^{1})\}\right].\]
Superscripts denote otherwise identical economies respectively excluding and including the event risk; the expectation integrates physical and economic shocks. Determine an identified or transparently bounded range for \(L\) using defensible transition probabilities and endogenous adaptation. Separate uncertainty about physical hazards from uncertainty about economic propagation, and test whether allowing migration, insurance, and correlated regional damages reverses policy rankings. This is a proposed quantitative model comparison for one event, rather than a claim that tipping points lack existing economic analyses.

\par\textbf{Formulation and scope.} Proposed scoped research specification. The original catalogue prompt was a synthesis, and no canonical conjecture or verified original formulation is claimed. The supporting source, when retained, establishes context or existing results rather than this exact open question.

\par\textbf{Open-question provenance.} An explicit published open-question statement was not verified. Sources below are supporting literature and must not be treated as a first listing.

\par\textbf{Historical priority.} No readable primary passage explicitly stating this precise open question was verified. A supporting paper, abstract, or topic citation is insufficient to establish origin or unresolved status.

\subsection{Short English statement}
Bound the century-long welfare cost of one specified climate tipping event when migration, investment, and adaptation respond, separating physical uncertainty from economic uncertainty.

\subsection{Sources}
\begin{itemize}\raggedright
\item[\textnormal{[S1]}]\hypertarget{source:30007713:1}{} Year not verified. A Guide to Macroeconomics and Climate Change (2025). Suggested supporting reference from the initial catalogue; exact open-question passage not verified..
\url{https://www.nber.org/papers/w33567}.
{\small\textit{Unverified bibliographic lead Bibliographic lead only. It is not evidence of an original explicit open-question statement.}}

\item[\textnormal{[S2]}]\hypertarget{source:30007713:2}{} Year not verified. Synthesis of evidence yields high social cost of carbon due to structural model variation and uncertainties (2024). Suggested supporting reference from the initial catalogue; exact open-question passage not verified..
\url{https://pmc.ncbi.nlm.nih.gov/articles/PMC11670083/}.
{\small\textit{Unverified bibliographic lead Bibliographic lead only. It is not evidence of an original explicit open-question statement.}}
\end{itemize}
\end{document}
